\chapter{Что вычисляют \nt{СЕРО}-нейроны}
\label{sec:serocomputer}

\section{Первый широковещательный канал}

До сих пор все нейроны книги говорили по адресной шине: \nt{ГЛУТ} и \nt{ГАМК} посылали импульсы определённым получателям, и мы выяснили, что такая сеть вычисляет и на каком языке она говорит. Теперь переходим ко второй шине из раздела~\ref{sec:buses} --- широковещательной. Первый её канал --- \nt{СЕРО}. Как и раньше, разрешим себе только то, что бывает в живом организме.

В этой главе появятся три устройства, которыми потом будут пользоваться все широковещательные каналы: нейрон, который при ровном фоновом притоке работает сам, пока его не остановили; провод, который на самом деле объём ткани; и медленная концентрация вместо отдельных импульсов. \nt{ДОФА}, \nt{НОРА} и \nt{АЦЕТ} в главах~\ref{sec:dofa}--\ref{sec:acet} будут собраны из тех же деталей.

С точки зрения инженера у \nt{СЕРО}-нейрона четыре важных отличия.

\emph{Первое: знак выбирает получатель.} У серотонина есть рецепторы обоих знаков, и правило~\eqref{eq:dalesign} здесь не действует: знак определяет не отправитель, а рецептор получателя (утверждение~\ref{prop:receptor})~\cite{BarnesSharp1999}.

\emph{Второе: \nt{СЕРО}-нейрон работает сам, если есть ровный фон.} В бодрствовании он ровно, несколько раз в секунду, выдаёт импульсы, и толчок на каждый импульс ему не нужен: это ритмоводитель. Но ему нужен постоянный фоновый приток. В срезе мозга, где этого притока нет, большинство таких клеток молчит; ровный ритм возвращается, если подать норадреналин через $\alpha_1$-рецепторы, а их блокада его гасит~\cite{VandermaelenAghajanian1983}. В организме этот фон приходит от \nt{НОРА}. Для схемы это генератор, которому нужно питание: пока оно есть, он работает сам, пока его не остановит торможение.

\emph{Третье: он медленный.} Почти все рецепторы серотонина метаботропные: они не открывают канал сами, а запускают внутри клетки цепочку реакций. Ответ медленный: тормозящий ток через главный рецептор нарастает и спадает в пределах примерно секунды~\cite{CourtneyFord2016} --- в десятки раз дольше, чем ответ быстрого синапса \nt{ГЛУТ}.

\emph{Четвёртое: он выбрасывает медиатор не в провод, а в объём.} В областях, куда идут выходы \nt{СЕРО}-нейронов, серотонин часто выходит не в узкую щель, а в межклеточное пространство и расплывается вокруг~\cite{Bunin1998}. Получатель чувствует концентрацию и не может знать, от какого отправителя пришла молекула.

Отдельные импульсы при таких временах уже неважны, поэтому описывать сеть будем частотами $r_j$ и концентрациями. Концентрация $c$ серотонина в объёме, куда выбрасывают медиатор нейроны $j$, растёт от выброса и убывает, потому что медиатор откачивают обратно:
\begin{empheq}[box=\keybox]{equation}
\tau_c\,\frac{\dd c}{\dd t} \;=\; -c \;+\; \kappa\sum_j r_j ,
\qquad
r_i \;=\; f_i\bigl(b_i + s_i\,g\,c\bigr), \quad s_i=\pm1 .
\label{eq:seroneuron}
\end{empheq}
Это ещё один способ прочитать поток импульсов, дополняющий утверждение~\ref{prop:reader}: объём не видит ни отдельных импульсов, ни их порядка, а только частоту, усреднённую за время $\tau_c$. Первое уравнение --- та же RC-цепь, что и мембрана: $\tau_c$ --- время жизни медиатора в объёме; прямо оно не измерено, и мы принимаем его порядка секунды, $\kappa$ --- сколько медиатора даёт один импульс. Второе описывает получателя: $b_i$ --- его собственный приток, у ритмоводителя положительный (в него входит и ровный фоновый вход от \nt{НОРА}); $g$ --- чувствительность рецептора; знак $s_i$ выбирает получатель; $f_i$ --- неубывающая передаточная характеристика.

\section{НЕ за один нейрон}

Возьмём ритмоводитель с тормозящим рецептором, $s=-1$. Пока серотонина вокруг нет, он работает на своём собственном притоке $b$. Когда серотонин появляется, приток уменьшается на $gc$, и если $gc>b$, нейрон замолкает. Это НЕ: выход есть, когда входа нет. На него ушёл один нейрон (рис.~\ref{fig:serogates}). Утверждение~\ref{prop:monotone} здесь неприменимо: оно требовало положительных весов.

Если на тот же ритмоводитель действует серотонин из двух разных объёмов, он молчит, когда серотонин есть хотя бы в одном. Это ИЛИ-НЕ, а из ИЛИ-НЕ можно собрать любую логическую функцию. Двухпроводный код \nt{ГЛУТ}-сети здесь не нужен: отсутствие сигнала вычисляют нейроны, которые работают, пока их не остановили.

\begin{figure}[t]
\centering
\begin{tikzpicture}[>=Stealth,
  nrn/.style={circle,draw,very thick,fill=blue!6,minimum size=0.95cm,inner sep=0pt},
  pace/.style={nrn,double,double distance=1.2pt},
  inh/.style={-{Bar[width=7pt]},thick,red!70!black}]
\node[pace] (N1) at (0,0) {};
\draw[inh] (-1.6,0) node[left,black] {$c$} -- (N1);
\draw[->,very thick,red!70!black] (N1) -- (1.3,0) node[right,black] {$\bar c$};
\node[below=0.55cm] at (0,0) {\small НЕ};
\node[pace] (N2) at (5,0) {};
\draw[inh] (3.4,0.45) node[left,black] {$c_1$} -- (N2);
\draw[inh] (3.4,-0.45) node[left,black] {$c_2$} -- (N2);
\draw[->,very thick,red!70!black] (N2) -- (6.3,0) node[right,black] {$\overline{c_1\vee c_2}$};
\node[below=0.55cm] at (5,0) {\small ИЛИ-НЕ};
\end{tikzpicture}
\caption{Логика из \nt{СЕРО}-нейронов. Двойной кружок --- ритмоводитель: при ровном фоновом притоке он работает сам, пока его не затормозят. Линия с чертой --- тормозящий рецептор, $s=-1$; $c$, $c_1$, $c_2$ --- концентрации серотонина в объёмах вокруг получателя. Слева --- НЕ, справа --- ИЛИ-НЕ.}
\label{fig:serogates}
\end{figure}

\begin{keyblock}
\begin{proposition}[Полнота \nt{СЕРО}-логики]
\label{prop:serocomplete}
Если в сети есть ритмоводители с тормозящими рецепторами и выбросы в разные объёмы не смешиваются, сеть~\eqref{eq:seroneuron} может вычислить любую логическую функцию, и каждый бит передаётся одним проводом.
\end{proposition}
\end{keyblock}

Логически \nt{СЕРО}-сеть сильнее \nt{ГЛУТ}, но условие «выбросы в разные объёмы не смешиваются» в мозге не выполняется (раздел~\ref{sec:volume}).

\section{Отрицательная обратная связь}

У серотониновых нейронов есть тормозящие рецепторы на них самих~\cite{Sprouse1987}. Серотонин, который они выбросили, возвращается к ним же и притормаживает их. Для инженера это знакомая схема: усилитель с отрицательной обратной связью (рис.~\ref{fig:serofeedback}).

Что здесь измерено, а что --- модель. В самом ядре шва, где сидят \nt{СЕРО}-нейроны, серотонин тормозит соседей не через общий объём, а точка-точка, через синапсы: переносчик быстро убирает его и не даёт накопиться вне щели. Одиночный выброс даёт лишь короткую паузу в разряде, а средняя частота от этого не меняется~\cite{CourtneyFord2016}. Поэтому петля ниже, замкнутая через общий объём, --- модель: мы считаем, что делала бы такая обратная связь, если бы она работала на уровне средних частот.

\begin{figure}[t]
\centering
\begin{tikzpicture}[>=Stealth,
  nrn/.style={circle,draw,very thick,fill=blue!6,minimum size=0.95cm,inner sep=0pt},
  pace/.style={nrn,double,double distance=1.2pt},
  blk/.style={draw,very thick,rounded corners=2pt,fill=orange!10,minimum height=0.9cm,minimum width=2.2cm,align=center,font=\small},
  inh/.style={-{Bar[width=7pt]},thick,red!70!black}]
\node[pace] (N) at (0,0) {};
\node[blk] (V) at (3.6,0) {объём\\$\tau_c$};
\draw[->,very thick] (N) -- node[above] {\small $r$} (V);
\draw[->,very thick] (V) -- (6.0,0) node[right] {\small $c$ --- всем получателям};
\draw[inh] (V.south) -- ++(0,-0.9) -| node[pos=0.25,below] {\small $-g\,c$} (N.south);
\draw[->,thick,blue!60!black] (-1.7,0) node[left,black] {\small $b$} -- (N);
\end{tikzpicture}
\caption{\nt{СЕРО}-нейрон с обратной связью через собственный медиатор: концентрация $c$ идёт ко всем получателям и тормозит сам нейрон. В модели это регулятор уровня; в мозге такая роль петли --- гипотеза.}
\label{fig:serofeedback}
\end{figure}

Посчитаем, что делает эта петля. Пусть передаточная характеристика линейна, $f(u)=au$ при $u>0$. В равновесии $c=\kappa r$ и $r=a(b-gc)$. Подставив одно в другое, получим
\begin{empheq}[box=\keybox]{equation}
r \;=\; \frac{a\,b}{1+G}, \qquad G \;=\; a\,g\,\kappa .
\label{eq:feedback}
\end{empheq}
Величина $G$ --- усиление петли. Это та же формула, что и для любого усилителя с отрицательной обратной связью, и она даёт те же два выигрыша.

\emph{Устойчивость к разбросу.} Пусть возбудимость нейрона $a$ изменилась на долю $\varepsilon$. Без обратной связи на ту же долю изменилась бы и частота. С обратной связью, при $G\gg1$, частота $r\approx b/(g\kappa)$ от $a$ почти не зависит: её изменение в $1+G$ раз меньше. При $G=9$ разброс подавляется вдесятеро.

\emph{Быстрота.} Подставим $r=a(b-gc)$ в первое уравнение~\eqref{eq:seroneuron}: получим $\tau_c\,\dd c/\dd t=-(1+G)\,c+\kappa ab$. Концентрация подходит к своему уровню с постоянной времени $\tau_c/(1+G)$ --- в $1+G$ раз быстрее, чем без обратной связи.

Вот первое вычисление, которое могла бы выполнять \nt{СЕРО}-система: держать общий уровень серотонина в мозге на заданной отметке, как термостат держит температуру. Это гипотеза: в опыте автоторможение пока видно только как короткие паузы, не меняющие средней частоты.

\section{Из импульсов --- в уровень}

Второе вычисление спрятано в первом уравнении~\eqref{eq:seroneuron}. Нейроны выдают отдельные импульсы, а получатель чувствует концентрацию, которая сглаживает их за время $\tau_c$. Если частота источника меняется, концентрация следует за ней с запаздыванием:
\begin{equation}
c(t) \;=\; \frac{\kappa}{\tau_c}\int_{-\infty}^{t} e^{-(t-t')/\tau_c}\,\sum_j r_j(t')\,\dd t' .
\label{eq:lowpass}
\end{equation}
Это скользящее среднее активности источников за последние несколько $\tau_c$ --- фильтр нижних частот (рис.~\ref{fig:lowpass}). Так простейший цифро-аналоговый преобразователь пропускает импульсы через RC-цепь и превращает их частоту в уровень. \nt{СЕРО}-система делает то же с сотнями тысяч нейронов сразу.

Эта величина заодно и память. После того как источники замолчали, концентрация держится ещё время порядка $\tau_c$. Это не бит, а плавно тающий след.

\begin{figure}[t]
\centering
\begin{tikzpicture}[>=Stealth,x=1.45cm,y=1cm]
\draw[gray!70!black] (0,3.2) -- (8.1,3.2);
\node[left,font=\small] at (-0.05,3.4) {импульсы};
\node[above,font=\scriptsize,gray!40!black] at (1,3.65) {$2$ в секунду};
\node[above,font=\scriptsize,gray!40!black] at (3.5,3.65) {$8$ в секунду};
\node[above,font=\scriptsize,gray!40!black] at (6.5,3.65) {$2$ в секунду};
\draw[->,gray!70!black] (0,0) -- (8.3,0) node[right,black] {\small $t$};
\draw[->,gray!70!black] (0,0) -- (0,2.75) node[left,black] {\small $c$};
\foreach \s in {0.136,0.529,0.760,1.223,1.714,1.748,1.755,2.663,2.701,2.734,3.414,3.493,3.719,3.800,3.928,3.948,4.074,4.327,4.420,4.589,4.728,4.736,4.914,5.025,5.205,5.220,6.224,6.544,7.178}{\draw[thick,blue!60!black] (\s,3.2) -- (\s,3.6);}
\foreach \a/\b/\v in {0.000/0.136/0.000,0.136/0.529/1.000,0.529/0.760/1.675,0.760/1.223/2.330,1.223/1.714/2.466,1.714/1.748/2.509,1.748/1.755/3.425,1.755/2.663/4.402,2.663/2.701/2.775,2.701/2.734/3.672,2.734/3.414/4.553,3.414/3.493/3.306,3.493/3.719/4.055,3.719/3.800/4.235,3.800/3.928/4.905,3.928/3.948/5.316,3.948/4.074/6.211,4.074/4.327/6.476,4.327/4.420/6.028,4.420/4.589/6.493,4.589/4.728/6.483,4.728/4.736/6.642,4.736/4.914/7.589,4.914/5.025/7.351,5.025/5.205/7.579,5.205/5.220/7.331,5.220/6.224/8.221,6.224/6.544/4.012,6.544/7.178/3.914,7.178/8.000/3.076}{\draw[very thick,orange!85!black] plot[domain=\a:\b,samples=10] (\x,{0.25*\v*exp(-(\x-\a))});}
\foreach \s/\p/\q in {0.136/0.000/1.000,0.529/0.675/1.675,0.760/1.330/2.330,1.223/1.466/2.466,1.714/1.509/2.509,1.748/2.425/3.425,1.755/3.402/4.402,2.663/1.775/2.775,2.701/2.672/3.672,2.734/3.553/4.553,3.414/2.306/3.306,3.493/3.055/4.055,3.719/3.235/4.235,3.800/3.905/4.905,3.928/4.316/5.316,3.948/5.211/6.211,4.074/5.476/6.476,4.327/5.028/6.028,4.420/5.493/6.493,4.589/5.483/6.483,4.728/5.642/6.642,4.736/6.589/7.589,4.914/6.351/7.351,5.025/6.579/7.579,5.205/6.331/7.331,5.220/7.221/8.221,6.224/3.012/4.012,6.544/2.914/3.914,7.178/2.076/3.076}{\draw[very thick,orange!85!black] (\s,0.25*\p) -- (\s,0.25*\q);}
\draw[thick,densely dashed,black] plot[domain=0:2,samples=30] (\x,{0.25*2*(1-exp(-\x))})
  plot[domain=2:5,samples=40] (\x,{0.25*(8-6.271*exp(-(\x-2)))})
  plot[domain=5:8,samples=40] (\x,{0.25*(2+5.688*exp(-(\x-5)))});
\foreach \x in {0,2,5,8}{\draw[gray!60!black] (\x,-0.05) -- (\x,0.05) node[below=3pt,font=\scriptsize] {\x\,с};}
\end{tikzpicture}
\caption{Из импульсов --- в уровень. Сверху --- импульсы \nt{СЕРО}-нейронов: две секунды редко, три секунды часто, потом снова редко. Снизу --- концентрация серотонина~\eqref{eq:lowpass} при $\tau_c=1\,\text{с}$. Штриховая линия --- то же, если бы импульсы шли идеально ровно.}
\label{fig:lowpass}
\end{figure}

\section{Провод --- это объём}
\label{sec:volume}

Вернёмся к условию утверждения~\ref{prop:serocomplete}. Серотонин, выброшенный поблизости разными отправителями, складывается в одну концентрацию.

\emph{Провод здесь --- не волокно, а объём ткани.} Два сигнала остаются разными, только если их выбрасывают в области, разнесённые дальше, чем успевает расплыться медиатор. За время $\tau$ молекула с коэффициентом диффузии $D$ уходит на
\begin{empheq}[box=\keybox]{equation}
\ell \;\sim\; \sqrt{D\,\tau}\, .
\label{eq:diffusionlength}
\end{empheq}
Извилистость межклеточных щелей замедляет диффузию небольшой молекулы примерно в $2.6$ раза~\cite{Sykova2008}, а коэффициент диффузии самого серотонина в ткани не измерен; по размеру молекулы оценим его в несколько единиц на $10^{-6}$~см$^2$/с. Если сигнал живёт $\tau\sim1$~с (тоже оценка), то $\ell\sim\sqrt{3\cdot10^{-6}}$~см $\approx20$~мкм. Адрес в такой сети --- это \emph{место}: получатель узнаёт, от кого сигнал, только по тому, где он сам находится.

Настоящие \nt{СЕРО}-нейроны устроены так, что этих отдельных адресов почти не остаётся. Выход каждого из них раскидан по огромным участкам мозга: ветви одной клетки доходят до многих далёких друг от друга областей~\cite{GagnonParent2014}. Мест выброса на клетку по оценке около $10^5$ (таблица~\ref{tab:types}); прямо их не подсчитали. «Провода» разных \nt{СЕРО}-нейронов перекрываются и смешиваются. Совсем в один провод они всё же не сливаются: \nt{СЕРО}-нейроны делятся на несколько подсистем, которые шлют выход в разные области и по-разному отвечают на одно и то же событие~\cite{Ren2018}. Но таких подсистем единицы, а не сотни тысяч.

\begin{keyblock}
\begin{proposition}[Число независимых проводов]
\label{prop:volumewires}
В сети с объёмной передачей число независимых сигналов ограничено не числом нейронов, а числом непересекающихся областей размером $\ell\sim\sqrt{D\tau}$, в которые они выбрасывают медиатор. Если выход каждого нейрона раскидан по большому объёму, вся сеть передаёт лишь несколько общих переменных.
\end{proposition}
\end{keyblock}

Поэтому логические схемы утверждения~\ref{prop:serocomplete} в мозге собрать не из чего. А регулятору~\eqref{eq:feedback} и фильтру~\eqref{eq:lowpass} отдельные провода не нужны: им хватает одной общей величины. Фильтр прямо следует из измеренного выброса в объём в областях назначения~\cite{Bunin1998}; регулятор уровня --- гипотеза.

\section{Горизонт планирования}
\label{sec:horizon}

Для чего может пригодиться одна общая медленная величина? Хороший кандидат --- параметр, который должен быть одинаков для всей сети и меняться не быстрее, чем за секунды или минуты. В главе~\ref{sec:neurontypes} такой параметр уже появлялся: коэффициент $\gamma$ в формуле ошибки~\eqref{eq:rpe}, который говорит, насколько будущая награда дешевле немедленной. В непрерывном времени его место занимает \emph{горизонт} $\tau_\gamma$: награда $R$, до которой ждать время $D$, стоит сейчас $R\,e^{-D/\tau_\gamma}$.

Горизонт решает, стоит ли ждать. Пусть можно сразу взять маленькую награду $r_0$ или подождать большую $R$. Ждать выгодно, пока
\begin{empheq}[box=\keybox]{equation}
D \;<\; \tau_\gamma\,\ln\frac{R}{r_0} .
\label{eq:patience}
\end{empheq}
При $\tau_\gamma=2\,\text{с}$ и втрое большей отложенной награде стоит ждать до $2.2\,\text{с}$; если горизонт вдвое длиннее, --- до $4.4\,\text{с}$. Терпение пропорционально горизонту.

Формула~\eqref{eq:patience} опирается на экспоненциальное обесценивание. Измеренное обесценивание на задержках в секунды ближе к гиперболе $R/(1+D/\tau_\gamma)$: так у обезьян падают и ценность отложенной награды в выборе, и ответ \nt{ДОФА}-нейронов на её предвестник~\cite{KobayashiSchultz2008}. Для гиперболы условие~\eqref{eq:patience} заменяется на $D<\tau_\gamma\,(R/r_0-1)$: зависимость от отношения наград уже не логарифмическая, а линейная, и при тех же числах граница ожидания сдвигается почти вдвое. Гипербола получается из той же экспоненты, если задержка случайна (задача~\ref{ex:05:7}), а пропорциональность терпения масштабу времени сохраняется.

По гипотезе, горизонт и задаёт \nt{СЕРО}~\cite{Doya2002}. Для такой роли у него всё есть: один общий уровень на всю сеть, который меняется за секунды. Есть и прямые опыты: если искусственно включить серотониновые нейроны, животное дольше ждёт отложенную награду, прежде чем бросить ожидание~\cite{Miyazaki2014}, и дольше пытается добыть награду там, где она стала редкой~\cite{Lottem2018}. Как именно концентрация серотонина превращается в $\tau_\gamma$ внутри сети, неизвестно. В следующей главе горизонт войдёт в ошибку, которую рассылает \nt{ДОФА}.

\section{Итог}

\begin{table}[t]
\centering
\small
\begin{tabular}{>{\raggedright}p{3.6cm}>{\raggedright}p{4.3cm}>{\raggedright\arraybackslash}p{4.3cm}}
\toprule
& \nt{ГЛУТ} & \nt{СЕРО} \\
\midrule
время реакции & миллисекунды & доли секунды --- секунда \\
знак действия & только $+$ & $+$ и $-$, выбирает получатель \\
без входа & молчит & работает сам при ровном фоновом притоке \\
адресация & точка --- точка & объём, адрес --- место \\
НЕ & только через датчики «выключения» & один нейрон \\
память & самоподдерживающаяся активность, гаснет от усталости & концентрация, тает за время $\tau_c$ \\
главные вычисления & совпадения, задержки, логика, последовательности & сглаживание; регулирование уровня и горизонт $\tau_\gamma$ (гипотезы) \\
\bottomrule
\end{tabular}
\caption{Что вычисляют сети из одних \nt{ГЛУТ}-нейронов и из одних \nt{СЕРО}-нейронов.}
\label{tab:glutvssero}
\end{table}

Как логический элемент (таблица~\ref{tab:glutvssero}) \nt{СЕРО}-нейрон богаче \nt{ГЛУТ}, но как система связи он --- широковещательная рассылка: медленная и почти без адресов. Поэтому он не пытается быть процессором, а сглаживает и рассылает несколько общих величин, о которых шла речь в главе~\ref{sec:neurontypes}, --- по гипотезе, в том числе горизонт планирования. Ритмоводитель, объём и медленная концентрация встретятся нам ещё трижды: у \nt{ДОФА}, \nt{НОРА} и \nt{АЦЕТ}.

\section{Задачи}

\begin{exercise}[\lvlA]
\label{ex:05:1}
\emph{Провод --- это объём: оценка.} Примите для серотонина в ткани $D=3\cdot10^{-6}$~см$^2$/с (раздел~\ref{sec:volume}).
(а)~По~\eqref{eq:diffusionlength} найдите $\ell$ для сигналов, живущих $\tau=0.1\,\text{с}$, $1\,\text{с}$ и $1$~мин.
(б)~Сколько непересекающихся кубиков со стороной $\ell(1\,\text{с})$ помещается в $1\,\text{мм}^3$ ткани? Это верхняя граница числа независимых проводов утверждения~\ref{prop:volumewires} в таком объёме.
(в)~За какое время медиатор расплывётся диффузией на $1$~мм и на $5$~см? Чем тогда обеспечена широковещательность \nt{СЕРО} --- диффузией или ветвлением выходного провода с $\sim10^5$ мест выброса (таблица~\ref{tab:types})? Что остаётся на долю диффузии?
\end{exercise}

\begin{exercise}[\lvlA]
\label{ex:05:2}
\emph{Регулятор уровня: вывод.}
(а)~Из~\eqref{eq:seroneuron} при $f(u)=au$ выведите~\eqref{eq:feedback} и покажите, что относительная чувствительность $S=(a/r)\,\partial r/\partial a$ равна $1/(1+G)$ точно, а не только при $G\gg1$.
(б)~При $G=9$ возбудимость $a$ выросла на $20\,\%$. На сколько процентов изменится $r$ (без линеаризации)?
(в)~Пусть $f(u)=a\varphi(u)$ с произвольной гладкой неубывающей $\varphi$. Докажите, что в окрестности равновесия $S=1/(1+G)$ и постоянная времени равна $\tau_c/(1+G)$, где $G=a\,\varphi'(u^*)\,g\kappa$, $u^*=b-gc^*$. Где на сигмоиде $\varphi$ регулятор перестаёт работать?
\end{exercise}

\begin{exercise}[\lvlB]
\label{ex:05:3}
\emph{Медленный рецептор в петле.} Рецепторы \nt{СЕРО} метаботропные, и их ответ запаздывает.
(а)~Опишите рецептор RC-звеном: $\tau_r\,\dd u/\dd t=-u+b-gc$, $r=au$, а объём --- первым уравнением~\eqref{eq:seroneuron}. Выведите характеристическое уравнение замкнутой петли, её собственную частоту $\omega_n$ и коэффициент затухания $\zeta$. При $\tau_c=1\,\text{с}$ и $\tau_r=0.3\,\text{с}$ найдите $G$, при котором $\zeta=1/\sqrt2$, и перерегулирование переходного процесса при $G=9$.
(б)~Замените RC-звено чистой задержкой: $r(t)=a\bigl(b-gc(t-T)\bigr)$. Докажите, что при $G>1$ петля устойчива тогда и только тогда, когда
$T<T^*=\tau_c\arccos(-1/G)/\sqrt{G^2-1}$, а при $G\le1$ --- при любой задержке. Вычислите $T^*$ при $G=2$ и $G=9$.
(в)~Какое подавление разброса $1/(1+G)$ реально доступно регулятору, если задержка в петле --- сотни миллисекунд?
\end{exercise}

\begin{exercise}[\lvlB]
\label{ex:05:4}
\emph{Дробовой шум уровня.} Пусть импульсы всех источников образуют пуассоновский поток с частотой $\lambda$; по~\eqref{eq:lowpass} каждый импульс в момент $t_k$ добавляет к $c$ слагаемое $(\kappa/\tau_c)e^{-(t-t_k)/\tau_c}$.
(а)~Докажите, что $\langle c\rangle=\kappa\lambda$ и $\operatorname{Var}c=\kappa^2\lambda/(2\tau_c)$, так что коэффициент вариации $\mathrm{CV}=1/\sqrt{2\lambda\tau_c}$. (Разбейте время на интервалы $\dd t$, в каждом из которых импульс есть независимо с вероятностью $\lambda\,\dd t$.)
(б)~Найдите CV для $\lambda=2$ и $8$ импульсов в секунду при $\tau_c=1\,\text{с}$ (рис.~\ref{fig:lowpass}).
(в)~Для строго периодических импульсов найдите отношение максимума к минимуму установившейся «пилы» $c(t)$ при тех же частотах.
(г)~Каков CV, если все $3\cdot10^5$ \nt{СЕРО}-нейронов по $2$ импульса в секунду выбрасывают медиатор в один объём? Какое $\tau_c$ понадобилось бы одному нейрону с частотой $4$ импульса в секунду, чтобы получить $\mathrm{CV}=10\,\%$? Почему уровень делают многие нейроны, а не один медленный?
\end{exercise}

\begin{exercise}[\lvlB]
\label{ex:05:5}
\emph{Моделирование: обратная связь против шума.} Для этой главы отдельной модели в \texttt{models/} нет; напишите её на Python. $N=50$ ритмоводителей, $i$-й выдаёт пуассоновские импульсы с частотой $r_i=a_i[b-gc]_+$, где $a_i$ равномерно распределены в $[2.8,\,5.2]$ импульса в секунду; концентрация --- по~\eqref{eq:seroneuron}, каждый импульс увеличивает $c$ на $\kappa/\tau_c$, $\kappa=1/N$, $\tau_c=1\,\text{с}$. Интегрируйте по Эйлеру с шагом $1\ms$ в течение $300\,\text{с}$, первые $20\,\text{с}$ отбросьте. Прогоны: (i)~$b=1$, $g=2.25$ ($G=9$); (ii)~$b=0.1$, $g=0$ --- без обратной связи, тот же средний уровень; (iii)~оба варианта с $a_i$, умноженными на $1.2$. Измерьте средний уровень $c$, его CV, время, за которое автокорреляция $c$ падает в $e$ раз, и частоты отдельных нейронов.
(а)~Линеаризуйте систему вокруг равновесия и покажите, что с обратной связью $\mathrm{CV}^2=1/(2\tau_cN\bar ab)$: при том же среднем уровне обратная связь снижает CV в $\sqrt{1+G}$ раз. Предскажите все измеряемые величины.
(б)~Сравните модель с предсказанием.
(в)~Выравнивает ли обратная связь частоты отдельных нейронов? Объясните, что именно регулирует общий объём.
\end{exercise}

\begin{exercise}[\lvlB]
\label{ex:05:6}
\emph{Проектирование: XOR на \nt{СЕРО}-логике.} Вентиль --- один ритмоводитель с пороговой характеристикой: $r=r_0$, если $b-g\sum_kc_k>0$, и $r=0$ иначе; он выбрасывает медиатор в свой объём, $\tau_c\,\dd c/\dd t=-c+\kappa r$. Логическая единица --- $c_1=\kappa r_0$, ноль --- $c=0$.
(а)~Покажите, что вентиль вычисляет ИЛИ-НЕ своих входов и может управлять следующими вентилями с тем же $g$, если $G'=g\kappa r_0/b>1$. Почему один выход можно подать на сколько угодно входов?
(б)~Соберите XOR$(A,B)$ из пяти вентилей ИЛИ-НЕ. Сколько отдельных объёмов нужно и на каком расстоянии их надо держать при $\tau_c=1\,\text{с}$ (задача~\ref{ex:05:1})?
(в)~Вход вентиля пересёк порог в момент $0$. Найдите, через какое время выход пересечёт порог $b/g$ при переключении $0\to1$ и $1\to0$. Покажите, что сумма этих времён минимальна при $G'=2$.
(г)~Найдите время установления XOR на самом длинном пути при $\tau_c=1\,\text{с}$ для $G'=4$ и $G'=2$.
\end{exercise}

\begin{exercise}[\lvlB]
\label{ex:05:7}
\emph{Терпение при неизвестной задержке.}
(а)~Выведите~\eqref{eq:patience}.
(б)~Животное соглашается ждать втрое большую награду, если ожидание не длиннее $6\,\text{с}$. Какой горизонт $\tau_\gamma$ это означает? До скольких секунд оно будет ждать, если включение \nt{СЕРО}-нейронов удвоит горизонт?
(в)~Пусть задержка $D$ случайна и распределена экспоненциально со средним $\bar D$. Покажите, что ждать выгодно при $\bar D<\tau_\gamma\,(R/r_0-1)$, и сравните с фиксированной задержкой при $\tau_\gamma=2\,\text{с}$, $R=3r_0$.
(г)~Докажите, что при любом распределении задержки со средним $\bar D$ ценность ожидания не меньше, чем при фиксированной задержке $\bar D$. Что это значит для поведения там, где задержки непредсказуемы?
\end{exercise}

\begin{exercise}[\lvlC]
\label{ex:05:8}
\emph{Как концентрация может задавать горизонт.} В разделе~\ref{sec:horizon} открыт вопрос, как уровень серотонина превращается в $\tau_\gamma$. В~\eqref{eq:ctd} горизонт входит только слагаемым $-V/\tau_\gamma$. Пусть \nt{ДОФА}-нейрон получает оценку $V$ двумя путями: прямым возбуждающим с весом $k_1$ и тормозящим через \nt{ГАМК}-посредника, который сглаживает сигнал с постоянной времени $\tau_d$, с весом $k_2$.
(а)~Покажите, что для медленно меняющегося $V$ сумма путей равна $(k_1-k_2)V+k_2\tau_d\,\dd V/\dd t$, и найдите $k_1$, $k_2$, воспроизводящие $\dd V/\dd t-V/\tau_\gamma$, при $\tau_d=0.1\,\text{с}$, $\tau_\gamma=2\,\text{с}$.
(б)~Пусть \nt{СЕРО} умножает прямой вес на $m$: $k_1\to mk_1$. Найдите $\tau_\gamma(m)$ и чувствительность $\dd\ln\tau_\gamma/\dd\ln m$ при $m=1$. На сколько процентов надо изменить $m$, чтобы удвоить горизонт? При каком $m$ горизонт становится бесконечным?
(в)~Рассогласование весов на долю $\mu$ даёт относительную ошибку горизонта $\approx\mu\tau_\gamma/\tau_d$, а конечное $\tau_d$ --- ложную ошибку $\approx\tau_d/\tau_\gamma$ (задача~\ref{ex:06:6}). Покажите, что обе ошибки нельзя сделать меньше $\sqrt\mu$ одновременно. Что это значит, если веса совпадают с точностью $1\,\%$?
(г)~Открытая часть. Критик учится по тому $\delta$, который реально вычисляет схема. Объясните, почему рассогласование весов тогда не ошибка, а просто другой горизонт, и почему высокая чувствительность из~(б) может быть полезна. Предложите механизм, в котором $\tau_\gamma$ зависит от $c$ без тонкой настройки, и опыт, который отличит его от механизма~(б).
\end{exercise}
